\subsection{Degree-four invariants fix the orbit}\label{sec:fivecopy-orbit}
The moment certificate fixes invariant values; we must now recover the state. The invariants needed for this
step have degree at most four and are linear in the four-copy marginal of the five-copy moment. Combined with the tangent-line lemma this gives one
equality theorem for every family point with $1/32<x<1/12$ (Theorem~\ref{thm:fivecopy-family-equality}), applied below to the
purity and in Section~\ref{sec:seven-entropy-equality} to the entropy.

\begin{theorem}[Degree-four invariants identify the icosahedral family; proved, FO-1]
\label{thm:moment-orbit-rigidity}
Fix $y\in(1/32,1/12)$. Let $z$ be a unit pure seven-qubit state whose thirteen
product-$\mathrm{SU}(2)$, site-permutation and complex-conjugation invariant polynomials
\[
 A_1,\ C,\ D,\ \mathcal E_2,\ \mathcal E_3,\ \mathcal E_4,\ B_1,\ S_1,\ Q,\ R,\ F,\ L,\ \pi_3,
\]
all of bidegree at most $(4,4)$ and defined below, take their values at $\psi_y$.
Then $z$ is local-unitarily and site-permutation equivalent to $\psi_y$.
In particular, relabelling the distinguished site as $0$ and the ordinary sites as $1,\ldots,6$,
in a suitable local frame its one-body marginals are $I/2$ and
\[
 T_{0j}=v e_ze_z^{\mathsf T},\qquad T_{ij}=w e_ze_z^{\mathsf T}\quad(i,j>0),
 \qquad v=\frac{1-8y}{3},\quad w=\frac{32y-1}{15}.
\]
The interval contains both $37/524$ and the entropy parameter $x_*$.
\end{theorem}

\begin{proof}
Write $r_i^a=\langle\sigma_a^{(i)}\rangle$ and
$T_{ij}^{ab}=\langle\sigma_a^{(i)}\sigma_b^{(j)}\rangle$.
At site $i$, concatenate $T_{ij}$ for $j>i$ and $T_{ji}^{\mathsf T}$
for $j<i$ into the real $3\times18$ matrix $M_i$.  The incident-axis
sum of squares
\[
 C=\sum_i\sum_{\substack{a<b\\c<d}}
    (M_i^{ac}M_i^{bd}-M_i^{ad}M_i^{bc})^2
  =\frac12\sum_i\bigl[(\operatorname{tr}M_iM_i^{\mathsf T})^2
  -\operatorname{tr}(M_iM_i^{\mathsf T})^2\bigr]
\]
has bidegree $(4,4)$ in $z,\bar z$.  Also sum, over ordered quadruples
of distinct sites, the nonnegative tetrad squares
\[
 D_{ijkl}=\sum_{a,b,c,d}
 (T_{ij}^{ab}T_{kl}^{cd}-T_{ik}^{ac}T_{jl}^{bd})^2.
\]
The site-sum $D$ likewise has bidegree $(4,4)$.
Local $SU(2)$ actions rotate each Pauli index orthogonally; each
displayed squared tensor norm is unchanged.  Site permutations
relabel summands.  Complex conjugation reflects the $Y$ Pauli axis,
another orthogonal action.  Hence $C,D$ and
$A_1=\sum_{i,a}(r_i^a)^2$ are conjugation-even invariant polynomials.

By Theorem~\ref{thm:ico} and its proof, at $\psi_y$ every one-body marginal is $I/2$ and the two-body computational-basis marginals are $\rho_{si}=\operatorname{diag}(10u,5u+y,5u+y,10u)$ and $\rho_{ij}=\operatorname{diag}(\tfrac12-8u,8u,8u,\tfrac12-8u)$ with $u=(1-2y)/30$. In this displayed frame the axial pair tensors are $T_{0j}=ve_ze_z^{\mathsf T}$ and $T_{ij}=we_ze_z^{\mathsf T}$ with $v=2(10u)-2(5u+y)=(1-8y)/3$ and $w=2(\tfrac12-8u)-2(8u)=(32y-1)/15$, so $A_1=C=D=0$ at $\psi_y$ (checked independently by FO-1). On the stated interval $v>w>0$.
The hypothesis forces $A_1(z)=C(z)=D(z)=0$.
Thus all one-body Bloch vectors vanish and $\operatorname{rank}M_i\le1$
for every $i$.  At each site choose a unit vector $n_i$ spanning the
column space of $M_i$, or any unit vector if $M_i=0$.
Every $T_{ij}$ then has the form $c_{ij}n_i n_j^{\mathsf T}$.
Product local rotations send all $n_i$ to $e_z$, leaving a real
symmetric scalar edge array $(c_{ij})$.  The tetrad equalities become
\begin{equation}\label{eq:fivecopy-pair-tetrads}
 c_{ij}c_{kl}=c_{ik}c_{jl}=c_{il}c_{jk}
 \qquad(i,j,k,l\text{ distinct}).
\end{equation}

The following four scalar graph invariants are again conjugation-even,
site-symmetric and of bidegree at most $(4,4)$:
\[
 \mathcal E_2=\sum_{i<j}\|T_{ij}\|_F^2,\qquad
 \mathcal E_3=\sum_{i<j<k}\operatorname{tr}(T_{ij}T_{jk}T_{ki}),
\]
\[
 \mathcal E_4=\frac13\sum_{i<j<k<l}
     \sum_{\text{three cycles on }\{i,j,k,l\}}
       \operatorname{tr}(T_{ij}T_{jk}T_{kl}T_{li}),
 \qquad B_1=\sum_{i<j}\|T_{ij}\|_F^4.
\]
In the $\mathcal E_4$ expression the indices inside the trace follow each listed cycle.
Their family values are
\[
 \mathcal E_2=6v^2+15w^2,\quad \mathcal E_3=15v^2w+20w^3,\quad
 \mathcal E_4=20v^2w^2+15w^4,\quad B_1=6v^4+15w^4.
\]
In particular $\mathcal E_3>0$, so the edge array has a nonzero triangle.
By~\eqref{eq:fivecopy-pair-tetrads}, a vertex outside this triangle
is either connected to all its vertices or to none; in the latter
case it has no edge to any other vertex.  Two vertices connected to
the triangle have a nonzero mutual edge.  Thus the nonzero edges form
a complete graph on one subset, with the other vertices isolated.
For the scalar edge reconstruction, temporarily index the seven sites $1,\dots,7$, with a nonzero triangle on $1,2,3$, and put
$\tau=c_{12}c_{13}/c_{23}\ne0$ and
$\varepsilon=\operatorname{sgn}\tau$.
Set $t_1=\sqrt{|\tau|}$ and
$t_j=c_{1j}/(\varepsilon t_1)$ on the active subset, and $t_j=0$
elsewhere.  The tetrads give $c_{ij}=\varepsilon t_it_j$ for every pair.
Every nonzero triangle product therefore has sign $\varepsilon$;
since $\mathcal E_3>0$, $\varepsilon=1$.

Put $x_i=t_i^2\ge0$, $e_k=e_k(x_1,\ldots,x_7)$, and
$p_k=\sum_i x_i^k$.  The graph invariants above become
$\mathcal E_2=e_2$, $\mathcal E_3=e_3$, $\mathcal E_4=e_4$, $B_1=e_2(x_1^2,\ldots,x_7^2)$.
Newton's identities give
\[
 B_1=e_2^2-2e_1e_3+2e_4,\qquad
 e_1=\frac{\mathcal E_2^2+2\mathcal E_4-B_1}{2\mathcal E_3}=r+6\ell,
 \qquad r=\frac{v^2}{w},\quad\ell=w.
\]
They also determine $p_2=e_1^2-2e_2$ and $p_3=e_1^3-3e_1e_2+3e_3$.
For $\mu=e_1/7$ and $S_k=\sum_i(x_i-\mu)^k$, evaluating at the family profile $(r,\ell,\dots,\ell)$ with deviations $\frac67(r-\ell)$ and $-\frac17(r-\ell)$ (six times) gives (FO-1)
\[
 S_2=\frac67(r-\ell)^2>0,\qquad
 S_3=\frac{30}{49}(r-\ell)^3>0,\qquad
 S_3^2=\frac{25}{42}S_2^3.
\]
The sharp seven-variable skewness inequality is
$S_3\le(5/\sqrt{42})S_2^{3/2}$, with equality only for one
high value and six equal low values.  Indeed, maximising $S_3$ at
fixed zero mean and $S_2$ by Lagrange multipliers yields at most two
values; if the higher value occurs $m$ times, the ratio is
$(7-2m)/\sqrt{7m(7-m)}$, maximal only at $m=1$.
The recovered mean and variance give one value $r=v^2/w$ and six values
$\ell=w$. All are positive. Flipping the local axis directions makes all
$t_i$ positive, so there is one distinguished vertex, with incident weight
$\sqrt{r\ell}=v$, and every ordinary edge has weight $w$.

The pair correlations have now fixed the local axes and edge weights. Return to labels $0,\ldots,6$,
with $0$ distinguished. We next use the remaining invariants to recover the state in that frame.
The symbols $Q,R,F,L$ below denote seven-qubit invariants local to this proof. Let
$J_7=\epsilon^{\otimes7}K$ be the spin flip, with $K$ complex conjugation,
and form the $21$ one-body spin-flip vectors $V_{ia}=\sigma_a^{(i)}J_7z$
($0\le i\le6$, $a\in\{x,y,z\}$), their Gram matrix $G=V^\dagger V$,
the overlaps $h=V^\dagger z$ and $q_i=\sum_a|h_{ia}|^2$.
Put $d_i=\sum_{j\ne i}\|T_{ij}\|_F^2$ and define
\[
 S_1=\sum_iq_i,\qquad Q=\sum_iq_i^2,\qquad
 R=\sum_id_iq_i,
\]
\[
 F=\sum_{i\ne j}\bigl(\|T_{ij}\|_F^2q_i
                -\|T_{ij}^{\mathsf T}h_i\|^2\bigr),\qquad
 L=\sum_{i<j}\operatorname{Re}(h_i^\dagger T_{ij}h_j),
\]
where $T_{ij}=T_{ji}^{\mathsf T}$ for $i>j$.
Each summand of $F$ is nonnegative because the operator norm of a
matrix is at most its Frobenius norm.  The respective bidegrees in
$(z,\bar z)$ are $(2,2)$ for $S_1$, $(4,4)$ for $Q,R,F$, and $(3,3)$
for $L$.  All five quantities are invariant under product $SU(2)$
and site permutations: $J_7U=UJ_7$, so the Pauli indices of $h_i$
and $T_{ij}$ rotate together.  Under complex conjugation, with
$\Delta=\operatorname{diag}(1,-1,1)$, they transform as
$h_i(\bar z)=\Delta\overline{h_i(z)}$ and
$T_{ij}(\bar z)=\Delta T_{ij}(z)\Delta$; the real part makes $L$ invariant.
A global phase multiplies all $h_i$ by one phase.  Their values equal those of $\psi_y$ by hypothesis.

In the displayed frame of~\eqref{eq:ico}, $J_7|b\rangle=(-1)^{7-|b|}|\bar b\rangle$. The pair and complement strings have total weights $2$ and $5$, and the endpoint strings have weights $6$ and $1$, so their spin-flip signs must be retained separately. Comparing complementary amplitudes, the antilinear map $\sigma_z^{(s)}J_7$ acts on each string of the pair/complement pattern with sign $-1$ and on the endpoint pattern with sign $+1$. At an ordinary site $i$ the map is $\sigma_z^{(i)}J_7=(-1)^{b_s+b_i}\sigma_z^{(s)}J_7$, which reverses the sign on both endpoints and on the $10$ pair/complement strings with $b_i\ne b_s$; these $10$ strings, now with sign $+1$, and the other $20$, with sign $-1$, contribute $-10u$, while the endpoints contribute $-2y$. The transverse components of every $h_i$ vanish: $J_7\psi_y$ has the same support as $\psi_y$ ($s=0$ with ordinary weights $\{2,6\}$, $s=1$ with $\{0,4\}$), $X_i$ or $Y_i$ at an ordinary site changes the parity of the ordinary weight, and $X_s$ or $Y_s$ moves a string into the other branch, whose weight set is disjoint. This yields the family frame values (checked by FO-1)
\[
 h_s^*=-1+4y,\qquad h_o^*=-\frac{1+4y}{3},\qquad
 q_0^*=(h_s^*)^2,\quad q_o^*=(h_o^*)^2.
\]
Both $h_s^*$ and $h_o^*$ are negative. The pair profile gives
$d_0=6v^2$ and $d_i=v^2+5w^2$ for $i>0$, so $d_0>d_i$.
Since $F(\psi_y)=0$, the hypothesis gives $F(z)=0$; every pair weight
is positive, so each $h_i=\zeta_i e_z$ is axial.

The equalities for $S_1$ and $R$, with $d_0\ne d_i$, determine
$q_0=q_0^*$ and $\sum_{i=1}^6q_i=6q_o^*$.
The equality for $Q$ then gives
$\sum_{i=1}^6q_i^2=6(q_o^*)^2$.
Equality in Cauchy--Schwarz forces $q_i=q_o^*$ for every $i>0$.
Thus all $|\zeta_i|$ equal their nonzero family values.
Moreover
\[
 L(z)=\sum_{i<j}a_{ij}\operatorname{Re}(\bar\zeta_i\zeta_j)
 \le\sum_{i<j}a_{ij}|\zeta_i||\zeta_j|=L(\psi_y),
\]
where all $a_{ij}$ are positive.  The hypothesis on $L$ makes this an
equality, so every pair of phases agrees.  A global phase of $z$
can make $h_{0z}=h_s^*$ and $h_{iz}=h_o^*$ exactly.

The one-site Bloch vectors vanish; the Gram $G$ therefore has
identity diagonal $3\times3$ blocks and off-diagonal blocks $T_{ij}$.
Let $c$ have axial entries $c_{0z}=-1/2$, $c_{iz}=-1/4$ for $i>0$,
and zero transverse entries.  Since $T_{0j}=ve_ze_z^{\mathsf T}$ and $T_{ij}=we_ze_z^{\mathsf T}$, only axial entries of $Gc$ survive, with
$(Gc)_{0z}=-\tfrac12-\tfrac32v=h_s^*$ and $(Gc)_{iz}=-\tfrac14-\tfrac v2-\tfrac54w=h_o^*$; so $Gc=h$, and
$c^\dagger h=-\tfrac12h_s^*-\tfrac32h_o^*=1$.  Consequently
\[
 \|z-Vc\|^2=1-2\operatorname{Re}(c^\dagger h)
                   +c^\dagger Gc=0.
\]
Thus $z=(-Z_0/2-\sum_{i=1}^6Z_i/4)J_7z$.
Proposition~\ref{prop:kramers-conference-shell} therefore puts $z$ in the
complex pair/complement shell $z_{T,b}$, where $T$ is the shell's symmetric amplitude matrix (not the
correlation blocks $T_{ij}$ above), with entries $t_{ij}$. On that shell
$\langle Z_0Z_j\rangle=1-4b^2-4p_j$ and $\langle Z_iZ_j\rangle=1-4(p_i+p_j)+8|t_{ij}|^2$, with $\sigma$ and $p_j$ as
in~\eqref{eq:complexsos}. Equating these with $v$ and $w$, and using $\sum_jp_j=2\sigma=1-2b^2$, gives
$b^2=y$, $p_j=(1-2y)/6=\sigma/3$ and $|t_{ij}|^2=(1-2y)/30=\sigma/15$, so the second and third squares
of~\eqref{eq:complexsos} vanish and its first is $\pi_3(\psi_y)-121/917$. Since $z$ has the purity of $\psi_y$,
\eqref{eq:complexsos} gives $\sum_{i<j}|(TT^\dagger)_{ij}|^2=0$: $T$ is a scaled symmetric conference matrix.
Proposition~\ref{prop:confphase} classifies it by diagonal congruence and
ordinary-site permutations. Its displayed local unitaries preserve the
endpoint amplitudes for every $b$, so $z$ is equivalent to $\psi_y$.
At $y=37/524$ this recovers the purity orbit and the pair values
$19/131$ and $11/131$ checked independently by MO-1--MO-2.
\end{proof}

\begin{corollary}[Fourth-moment orbit rigidity]\label{cor:fourth-moment-orbit-rigidity}
Let $X_r^+(z)$ denote the normalised $r$-copy moment of a pure seven-qubit
state after product-local-unitary and site-permutation twirling and the
complex-conjugation average. Fix $y\in(1/32,1/12)$. If
$X_r^+(z)=X_r^+(\psi_y)$ for some $r\ge4$, then $z$ is local-unitarily and
site-permutation equivalent to $\psi_y$. In particular this applies to
both the purity and entropy optimisers.
\end{corollary}
\begin{proof}
A twirl-invariant polynomial of bidegree $(k,k)$ with $k\le r$ is a
linear functional of $X_r^+$, after multiplying by
$\operatorname{tr}(zz^*)^{r-k}=1$.  The thirteen invariants of
Theorem~\ref{thm:moment-orbit-rigidity} are of this kind with $k\le4$;
the $r$-copy moment determines the four-copy moment by partial trace.
\end{proof}

Lemma~\ref{lem:fivecopy-tangent-line} and Corollary~\ref{cor:fourth-moment-orbit-rigidity}
combine into the equality step shared by the purity and entropy theorems.

\begin{theorem}[Equality from a sharp five-copy identity on the family; proved]
\label{thm:fivecopy-family-equality}
Fix $y\in(1/32,1/12)$, and put $X_y=X_5^+(\psi_y)$ and
$X_y'=\frac{d}{dx}X_5^+(\psi_x)|_{x=y}$. Let $\Gamma_b$ be linear maps to
symmetric matrices, and $W_b$ matrices such that
the functionals $X\mapsto\operatorname{tr}[W_bSW_b^{\mathsf T}\Gamma_b(X)]$,
  with $b$ varying and $S$ ranging over symmetric matrices, span the critical space of
  linear functionals $f$ on the span of the twirled five-copy moments with
  $f(X_y)=f(X_y')=0$.
If a normalised pure seven-qubit state $z$ has
$W_b^{\mathsf T}\Gamma_b(X_5^+(z))W_b=0$ for every $b$, then
$X_5^+(z)=X_y$, and $z$ is local-unitarily and site-permutation
equivalent to $\psi_y$.
\end{theorem}
\begin{proof}
The family values above give $H(\psi_x)=155+96x$ for the invariant
$H=104S_1-25\mathcal E_2$, so $X_y'\ne0$.
Lemma~\ref{lem:fivecopy-tangent-line} gives $X_5^+(z)=X_y+tX_y'$.
Writing $h=H(\psi_y)$, both $H$ and $(H-h)^2$ are linear functionals
of the normalised five-copy moment, since their bidegrees are at most $(4,4)$.
The latter is critical, hence vanishes at $z$. But
$H(z)-h=96t$, so $t=0$. Thus $X_5^+(z)=X_y$, and
Corollary~\ref{cor:fourth-moment-orbit-rigidity} with $r=5$ gives the orbit.
\end{proof}

\begin{corollary}[seven-qubit purity minimum and rigidity; verified-arithmetic certificate, FP5-1]
\label{cor:seven-global-from-fivecopy}
Every normalised pure seven-qubit state satisfies $\pi_3\ge121/917$.
Equality holds precisely on the icosahedral local-unitary/permutation
orbit, and there the whole normalised five-copy moment equals that of
$\Psi_{\rm ico}$. Quantitatively, with $c>0$ the least eigenvalue of the
$C_b$,
\[
 \pi_3(z)-\frac{121}{917}\ \ge\ c\sum_b\operatorname{tr}
 \bigl[W_b^{\mathsf T}\Gamma_b(X_5^+(z))W_b\bigr].
\]
\end{corollary}
\begin{proof}
For a pure state the flipped five-copy moments are averages of rank-one
projectors, so every $\Gamma_b(X_5^+(z))\succeq0$. The identity of
Theorem~\ref{thm:fivecopy-sharp-certificate} gives the bound and the
displayed inequality. At equality every
$W_b^{\mathsf T}\Gamma_b(X_5^+(z))W_b$ vanishes, and
Theorem~\ref{thm:fivecopy-family-equality} with $y=37/524$ gives the
moment and orbit, using the spanning assertion of
Theorem~\ref{thm:fivecopy-sharp-certificate}. Conversely,
Theorem~\ref{thm:ico} gives equality on that orbit.
\end{proof}

\begin{corollary}[Labelled LU orbits in the family; proved, EO-1]
\label{cor:icosahedral-labelled-orbits}
For each fixed $y\in(1/32,1/12)$, the states obtained by permuting the
qubits of $\psi_y$ form exactly $42$ distinct orbits under local unitaries
and global phase. They form one orbit when qubit permutations are also
allowed. Representatives are indexed by a choice of distinguished qubit
and a permutation of the last three of the six ordered ordinary labels,
leaving the first three ordinary labels fixed.
\end{corollary}
\begin{proof}
The incident pair weights are $d_s=6v^2$ at the distinguished site and
$d_o=v^2+5w^2$ elsewhere, with $d_s>d_o$. These labelled quantities are
LU invariants, so different choices of distinguished site give seven
different LU classes of configurations.

Fix that site. Every canonical permuted state has the same nonzero
positive axial pair profile. A local unitary mapping two such states
must send each local $Z$ axis to itself or its negative. Positivity of
every pair coefficient forces the seven signs to agree. Thus all local
unitaries are diagonal, or all are anti-diagonal. A simultaneous bit flip
sends the real conference matrix $C$ to $-C$, which is already realised
by diagonal congruence $D=iI$ and the local unitaries of
Proposition~\ref{prop:confphase}.

For diagonal unitaries, equality of the two nonzero endpoint amplitudes
shows that the pair matrices are related by $C'=DCD$ with $D$ diagonal
unitary. Conversely every such congruence is implemented by the local
unitaries displayed in that proposition, preserving the endpoints.
When $C,C'$ have real sign entries, $d_id_j\in\{1,-1\}$ for every
$i\ne j$. Using three distinct indices gives
$D=\eta\operatorname{diag}(\epsilon_i)$, where
$\eta\in\{1,i\}$ and $\epsilon_i\in\{1,-1\}$, up to an irrelevant
common sign. The exact LU criterion for ordinary permutations is therefore
\[
 C'=\tau\operatorname{diag}(\epsilon_i)
                C\operatorname{diag}(\epsilon_i),\qquad \tau\in\{1,-1\}.
\]
Normalise the first ordinary row to all positive signs. Orthogonality to
that row says that every row of the five-vertex core has two negative
entries. Its negative-edge graph is therefore a pentagon. There are
$(5-1)!/2=12$ labelled pentagons, all attained by ordinary permutations
of $C$. In the normalised frame the remaining equivalence above is
identity or complementation: after choosing $\epsilon_0=1$, the first
row forces $\epsilon_i=\tau$ for every $i>0$. Complementation pairs
the twelve pentagons without fixed points, leaving six classes.

The six permutations fixing ordinary labels $0,1,2$ and permuting
$3,4,5$ represent all six classes. In the base pentagon
$1\!-\!3\!-\!2\!-\!4\!-\!5$ the vertices $1,2$ are not adjacent, and a
permutation fixing them keeps them non-adjacent, whereas they are adjacent
in every complement; and two such images coincide only if the permutations
differ by a dihedral symmetry fixing two vertices, that is, the identity.
EO-1 independently checks all $720$ ordinary permutations: the switching
and switching-complement automorphism counts are $60$ and $60$, and
the six classes each contain $120$ permutations. It also verifies the
explicit diagonal local gauges on both amplitude-pattern basis vectors. Combining these with the seven
distinguished-site choices gives $7\cdot6=42$ LU orbits.
\end{proof}
